\documentclass[a4paper,12pt]{article}

\usepackage[T2A]{fontenc}
\usepackage[utf8]{inputenc}
\usepackage[russian]{babel}

\usepackage{amsmath,amssymb,mathtools}
\usepackage{helvet}
\renewcommand{\familydefault}{\sfdefault}

\usepackage{icomma}
\usepackage{geometry}
\geometry{left=20mm,right=20mm,top=20mm,bottom=22mm}

\usepackage{microtype}
\usepackage{tikz}
\usetikzlibrary{calc,arrows.meta,positioning,shapes.geometric}

\usepackage{pgfplots}
\pgfplotsset{compat=1.18}

\usepackage{tabularx,booktabs}
\usepackage{enumitem}
\usepackage{hyperref}
\usepackage{parskip}
\usepackage{fancyhdr}
\usepackage{setspace}
\usepackage{xcolor}

\definecolor{accent}{HTML}{008080}
\definecolor{darkaccent}{HTML}{004D4D}
\definecolor{textgray}{HTML}{333333}
\definecolor{softgray}{HTML}{F3F5F5}

\hypersetup{
    colorlinks=true,
    linkcolor=darkaccent,
    urlcolor=darkaccent
}

\onehalfspacing
\setlength{\parindent}{0pt}
\setlist[itemize]{leftmargin=6mm,itemsep=2pt,topsep=3pt}
\setlist[enumerate]{leftmargin=8mm,itemsep=5pt,topsep=3pt}

\pagestyle{fancy}
\fancyhf{}
\fancyhead[L]{\small\color{textgray}Условия встречи двух тел}
\fancyhead[R]{\small\color{textgray}Володя Хачатурян}
\fancyfoot[C]{\small\thepage}
\renewcommand{\headrulewidth}{0.4pt}
\renewcommand{\headrule}{\hbox to\headwidth{\color{accent}\leaders\hrule height \headrulewidth\hfill}}

\newcommand{\sectionline}{%
    \vspace{-0.4em}
    {\color{accent}\rule{\linewidth}{0.7pt}}
    \vspace{0.4em}
}

\newcommand{\key}[1]{\textbf{\color{darkaccent}#1}}

\tikzset{
    flowstart/.style={
        ellipse,
        draw=darkaccent,
        line width=0.9pt,
        minimum width=35mm,
        minimum height=11mm,
        align=center
    },
    flowaction/.style={
        rectangle,
        rounded corners=2mm,
        draw=accent,
        line width=0.9pt,
        text width=62mm,
        minimum height=13mm,
        align=center,
        inner sep=3mm
    },
    flowdecision/.style={
        diamond,
        draw=darkaccent,
        line width=0.9pt,
        aspect=2.3,
        text width=34mm,
        align=center,
        inner sep=1.5mm
    },
    flowarrow/.style={
        -{Latex[length=3mm,width=2mm]},
        line width=0.9pt,
        draw=textgray
    }
}

\begin{document}

% ============================================================
% ТИТУЛЬНЫЙ ЛИСТ
% ============================================================

\begin{titlepage}
\thispagestyle{empty}

\begin{center}

\vspace*{23mm}

{\Large\color{darkaccent}\textbf{ФИЗИКА}}

\vspace{10mm}

{\Huge\bfseries Чек-лист}

\vspace{5mm}

{\LARGE\bfseries
Условия встречи двух тел\\[2mm]
и чтение графиков движения}

\vspace{14mm}

{\large
Равномерное прямолинейное движение\\
Координата, время, скорость, график \(x(t)\)
}

\vfill

{\large
\textbf{Володя Хачатурян}
}

\vspace{8mm}

{\large
Лёвин Артём Александрович\\
эксклюзивно для Володи Хачатурян
}

\vspace{8mm}

{\large \today}

\vfill

\begin{tikzpicture}[x=1cm,y=1cm]
\draw[
    accent,
    line width=1.1pt
]
(0,0)
.. controls (3.0,0.8) and (6.0,-0.8) ..
(9,0)
.. controls (11.5,0.6) and (13.2,0.2) ..
(15,0);
\end{tikzpicture}

\end{center}

\end{titlepage}

% ============================================================
% ОСНОВНОЙ ЧЕК-ЛИСТ
% ============================================================

\section*{\color{darkaccent}1. Главная идея занятия}
\sectionline

Для равномерного движения вдоль оси \(Ox\) координата тела определяется формулой
\[
\boxed{x=x_0+v_x t}.
\]

Здесь:
\[
x_0 \text{ --- начальная координата},\qquad
v_x \text{ --- проекция скорости на ось } Ox,\qquad
t \text{ --- время}.
\]

Для двух тел используются отдельные формулы:
\[
\boxed{x_1=x_{01}+v_{x1}t},
\qquad
\boxed{x_2=x_{02}+v_{x2}t}.
\]

\key{Физический смысл встречи:} в один и тот же момент времени оба тела находятся в одной точке пространства.

Поэтому основное условие встречи:
\[
\boxed{x_1=x_2}.
\]

Это центральная формула всего блока.

\section*{\color{darkaccent}2. Что означает знак скорости}
\sectionline

Знак \(v_x\) показывает направление движения относительно выбранной оси.

\begin{itemize}
    \item Если \(v_x>0\), координата с течением времени увеличивается.
    \item Если \(v_x<0\), координата с течением времени уменьшается.
    \item Если \(v_x=0\), координата остаётся постоянной.
\end{itemize}

Например:
\[
x=4+3t
\]
описывает движение в положительном направлении оси \(Ox\), поскольку \(v_x=3\).

Формула
\[
x=7-2t
\]
соответствует движению в отрицательном направлении, поскольку
\[
v_x=-2.
\]

\section*{\color{darkaccent}3. Встреча двух тел: формулы уже даны}
\sectionline

Пусть
\[
x_1=-2+t,
\qquad
x_2=5-t.
\]

В момент встречи:
\[
x_1=x_2.
\]

Следовательно,
\[
-2+t=5-t.
\]

Переносим слагаемые:
\[
t+t=5+2,
\]
\[
2t=7,
\]
\[
\boxed{t_{\text{встр}}=3{,}5\text{ с}}.
\]

Теперь подставляем найденное время в любую формулу координаты:
\[
x_{\text{встр}}=-2+3{,}5=1{,}5.
\]

Проверка по второму телу:
\[
x_{\text{встр}}=5-3{,}5=1{,}5.
\]

Итак,
\[
\boxed{t_{\text{встр}}=3{,}5\text{ с}},
\qquad
\boxed{x_{\text{встр}}=1{,}5\text{ м}}.
\]

\subsection*{Мини-проверка}

После нахождения времени встречи полезно вычислить обе координаты:
\[
x_1(t_{\text{встр}})
\quad\text{и}\quad
x_2(t_{\text{встр}}).
\]

Если значения совпали, условие встречи выполнено.

% ============================================================
% АЛГОРИТМ
% ============================================================

\clearpage
\thispagestyle{empty}

\begin{center}
{\LARGE\bfseries\color{darkaccent}
Алгоритм: как найти время и координату встречи}
\end{center}

\vspace{8mm}

\begin{center}
\begin{tikzpicture}[node distance=10mm]

\node[flowstart] (start) {Начало};

\node[flowaction, below=of start] (write)
{Запишите законы движения:
\[
x_1=x_{01}+v_{x1}t,\qquad
x_2=x_{02}+v_{x2}t
\]};

\node[flowaction, below=of write] (equal)
{Используйте условие встречи:
\[
x_1=x_2
\]};

\node[flowaction, below=of equal] (solve)
{Решите полученное уравнение относительно \(t\)};

\node[flowdecision, below=13mm of solve] (checktime)
{Получено ли физически допустимое время?};

\node[flowaction, below=13mm of checktime] (coord)
{Подставьте \(t_{\text{встр}}\) в одну из формул и найдите
\[
x_{\text{встр}}
\]};

\node[flowaction, below=of coord] (verify)
{Проверьте:
\[
x_1(t_{\text{встр}})
=
x_2(t_{\text{встр}})
\]};

\node[flowstart, below=of verify] (finish)
{Ответ};

\node[flowaction, right=25mm of checktime, text width=45mm] (reject)
{Проверьте исходные данные и физический смысл результата};

\draw[flowarrow] (start) -- (write);
\draw[flowarrow] (write) -- (equal);
\draw[flowarrow] (equal) -- (solve);
\draw[flowarrow] (solve) -- (checktime);
\draw[flowarrow] (checktime) -- node[right,font=\small]{да} (coord);
\draw[flowarrow] (coord) -- (verify);
\draw[flowarrow] (verify) -- (finish);

\draw[flowarrow]
(checktime.east)
-- node[above,font=\small]{нет}
(reject.west);

\draw[flowarrow]
(reject.north)
|- ([xshift=15mm]solve.east)
-- (solve.east);

\end{tikzpicture}
\end{center}

\clearpage

% ============================================================
% ДАННЫЕ В ЧИСЛОВОМ ВИДЕ
% ============================================================

\section*{\color{darkaccent}4. Формула движения в условии отсутствует}
\sectionline

Иногда в задаче указываются только начальные координаты и скорости.

Например:
\[
x_{01}=10\text{ м},
\qquad
v_{x1}=-2\text{ м/с},
\]
\[
x_{02}=-10\text{ м},
\qquad
v_{x2}=2\text{ м/с}.
\]

Сначала собираем законы движения:
\[
x_1=10-2t,
\]
\[
x_2=-10+2t.
\]

В момент встречи:
\[
10-2t=-10+2t.
\]

Получаем:
\[
20=4t,
\]
\[
\boxed{t_{\text{встр}}=5\text{ с}}.
\]

Координата встречи:
\[
x_{\text{встр}}=10-2\cdot5=0.
\]

Ответ:
\[
\boxed{t_{\text{встр}}=5\text{ с}},
\qquad
\boxed{x_{\text{встр}}=0}.
\]

\subsection*{Что требуется помнить}

При составлении формулы
\[
x=x_0+v_xt
\]
знак скорости сохраняется.

Если
\[
v_x=-2\text{ м/с},
\]
то
\[
x=x_0-2t.
\]

Если
\[
v_x=2\text{ м/с},
\]
то
\[
x=x_0+2t.
\]

% ============================================================
% ГРАФИК
% ============================================================

\section*{\color{darkaccent}5. Как увидеть встречу на графике}
\sectionline

На графике \(x(t)\):

\begin{itemize}
    \item по горизонтальной оси откладывается время \(t\);
    \item по вертикальной оси откладывается координата \(x\);
    \item точка пересечения двух графиков соответствует встрече тел.
\end{itemize}

Для примера рассмотрим
\[
x_1=5-t,
\qquad
x_2=1+\frac13t.
\]

\begin{center}
\begin{tikzpicture}
\begin{axis}[
    width=0.83\linewidth,
    height=8.3cm,
    axis lines=middle,
    xmin=0,
    xmax=6.5,
    ymin=0,
    ymax=6,
    xlabel={$t$, с},
    ylabel={$x$, м},
    xtick={0,1,2,3,4,5,6},
    ytick={0,1,2,3,4,5,6},
    grid=both,
    major grid style={line width=0.3pt},
    minor tick num=0,
    clip=false,
    samples=200,
    legend style={
        at={(0.97,0.97)},
        anchor=north east,
        draw=none,
        fill=none
    }
]
\addplot[
    domain=0:5,
    thick,
    darkaccent
]{5-x};
\addlegendentry{\(x_1=5-t\)}

\addplot[
    domain=0:6,
    thick,
    accent
]{1+x/3};
\addlegendentry{\(x_2=1+\frac13t\)}

\addplot[
    only marks,
    mark=*,
    mark size=2.4pt,
    textgray
]
coordinates {(3,2)};

\node[
    anchor=south west,
    font=\small
]
at (axis cs:3.05,2.08)
{\((3;\,2)\)};

\end{axis}
\end{tikzpicture}
\end{center}

Из графика видно:
\[
\boxed{t_{\text{встр}}=3\text{ с}},
\qquad
\boxed{x_{\text{встр}}=2\text{ м}}.
\]

Проверим аналитически:
\[
5-t=1+\frac13t.
\]

Тогда
\[
4=\frac43t,
\]
\[
t=3.
\]

Далее
\[
x=5-3=2.
\]

% ============================================================
% СЧИТЫВАНИЕ СКОРОСТИ
% ============================================================

\section*{\color{darkaccent}6. Как считать скорость с графика}
\sectionline

Для равномерного движения график \(x(t)\) является прямой.

Скорость показывает, насколько изменяется координата за единицу времени:
\[
\boxed{v_x=\frac{\Delta x}{\Delta t}}.
\]

Выбираются две удобные точки графика:
\[
(t_1;x_1),
\qquad
(t_2;x_2).
\]

Тогда
\[
\boxed{
v_x=
\frac{x_2-x_1}{t_2-t_1}
}.
\]

\subsection*{Пример 1. Убывающий график}

Пусть график проходит через точки
\[
(0;5)
\quad\text{и}\quad
(1;4).
\]

Тогда
\[
v_x=\frac{4-5}{1-0}=-1\text{ м/с}.
\]

Начальная координата:
\[
x_0=5\text{ м}.
\]

Следовательно,
\[
\boxed{x=5-t}.
\]

\subsection*{Пример 2. Возрастающий график}

Пусть график проходит через точки
\[
(0;1)
\quad\text{и}\quad
(6;3).
\]

Тогда
\[
v_x=\frac{3-1}{6-0}
=
\frac26
=
\frac13\text{ м/с}.
\]

Поэтому
\[
\boxed{x=1+\frac13t}.
\]

\subsection*{Быстрый визуальный контроль}

\[
\text{график возрастает}
\Longrightarrow
v_x>0,
\]
\[
\text{график убывает}
\Longrightarrow
v_x<0.
\]

Этот признак помогает проверить знак полученной скорости.

% ============================================================
% АЛГОРИТМ ГРАФИКА
% ============================================================

\clearpage
\thispagestyle{empty}

\begin{center}
{\LARGE\bfseries\color{darkaccent}
Алгоритм: восстановление закона движения по графику}
\end{center}

\vspace{8mm}

\begin{center}
\begin{tikzpicture}[node distance=10mm]

\node[flowstart] (start) {Начало};

\node[flowaction, below=of start] (xzero)
{Считайте значение \(x\) при \(t=0\):
\[
x_0=x(0)
\]};

\node[flowaction, below=of xzero] (points)
{Выберите две удобные точки графика:
\[
(t_1;x_1),\quad(t_2;x_2)
\]};

\node[flowaction, below=of points] (speed)
{Вычислите скорость:
\[
v_x=\frac{x_2-x_1}{t_2-t_1}
\]};

\node[flowdecision, below=13mm of speed] (sign)
{Совпадает ли знак \(v_x\) с направлением графика?};

\node[flowaction, below=13mm of sign] (formula)
{Составьте закон движения:
\[
x=x_0+v_xt
\]};

\node[flowaction, below=of formula] (control)
{Подставьте координаты одной узловой точки и выполните проверку};

\node[flowstart, below=of control] (finish)
{Формула восстановлена};

\node[flowaction, right=24mm of sign, text width=45mm] (recalc)
{Проверьте порядок вычитания и знаки при вычислении \(\Delta x\) и \(\Delta t\)};

\draw[flowarrow] (start) -- (xzero);
\draw[flowarrow] (xzero) -- (points);
\draw[flowarrow] (points) -- (speed);
\draw[flowarrow] (speed) -- (sign);
\draw[flowarrow] (sign) -- node[right,font=\small]{да} (formula);
\draw[flowarrow] (formula) -- (control);
\draw[flowarrow] (control) -- (finish);

\draw[flowarrow]
(sign.east)
-- node[above,font=\small]{нет}
(recalc.west);

\draw[flowarrow]
(recalc.north)
|- ([xshift=15mm]speed.east)
-- (speed.east);

\end{tikzpicture}
\end{center}

\clearpage

% ============================================================
% АБСЦИССА И ОРДИНАТА
% ============================================================

\section*{\color{darkaccent}7. Абсцисса и ордината на графике \(x(t)\)}
\sectionline

Для точки графика
\[
(t;x)
\]
первая координата соответствует горизонтальной оси, вторая --- вертикальной.

\begin{center}
\renewcommand{\arraystretch}{1.35}
\begin{tabularx}{0.9\linewidth}{@{}lXl@{}}
\toprule
\textbf{Термин} & \textbf{Что означает} & \textbf{На графике \(x(t)\)} \\
\midrule
Абсцисса &
координата точки по горизонтальной оси &
\(t\) \\
Ордината &
координата точки по вертикальной оси &
\(x\) \\
\bottomrule
\end{tabularx}
\end{center}

Например, точка
\[
(6;3)
\]
означает:
\[
t=6\text{ с},
\qquad
x=3\text{ м}.
\]

Эти значения можно подставить в
\[
x=x_0+v_xt.
\]

% ============================================================
% ОБРАТНАЯ ЗАДАЧА
% ============================================================

\section*{\color{darkaccent}8. Обратная задача: время встречи уже известно}
\sectionline

Пусть для первого тела
\[
x_1=3+t.
\]

Известно, что тела встретились через
\[
t_{\text{встр}}=5\text{ с}.
\]

Для второго тела
\[
x_2=x_{02}+2t,
\]
причём требуется найти \(x_{02}\).

Сначала определяем общую координату встречи по первому телу:
\[
x_{\text{встр}}
=
x_1(5)
=
3+5
=
8\text{ м}.
\]

В точке встречи координата второго тела также равна \(8\) м:
\[
8=x_{02}+2\cdot5.
\]

Следовательно,
\[
8=x_{02}+10,
\]
\[
\boxed{x_{02}=-2\text{ м}}.
\]

Итак,
\[
\boxed{x_{\text{встр}}=8\text{ м}},
\qquad
\boxed{x_{02}=-2\text{ м}}.
\]

\key{Ключевая идея:} координата встречи является общей для обоих тел.

% ============================================================
% ТИПИЧНЫЕ ОШИБКИ
% ============================================================

\section*{\color{darkaccent}9. Типичные ошибки}
\sectionline

\begin{enumerate}
    \item \textbf{Ошибка при переносе слагаемых.}

    В уравнении
    \[
    2t-3=t+2
    \]
    при переносе знак слагаемого меняется:
    \[
    2t-t=2+3.
    \]

    Отсюда
    \[
    t=5.
    \]

    \item \textbf{Потеря отрицательного знака начальной координаты.}

    Если
    \[
    x_0=-10,
    \]
    формула должна содержать именно \(-10\).

    \item \textbf{Путаница между временем и координатой встречи.}

    Сначала обычно находится
    \[
    t_{\text{встр}},
    \]
    затем
    \[
    x_{\text{встр}}.
    \]

    \item \textbf{Неверное решение уравнения \(0=4t\).}

    Правильный вывод:
    \[
    t=0.
    \]

    \item \textbf{Потеря знака скорости при чтении графика.}

    Убывающий график соответствует
    \[
    v_x<0.
    \]

    Возрастающий график соответствует
    \[
    v_x>0.
    \]

    \item \textbf{Перепутаны координаты узловой точки.}

    На графике \(x(t)\) запись
    \[
    (6;3)
    \]
    означает
    \[
    t=6,\qquad x=3.
    \]

    \item \textbf{Замазывание ошибочной записи.}

    Для аккуратного оформления достаточно зачеркнуть ошибочную запись одной линией и продолжить решение рядом.
\end{enumerate}

% ============================================================
% ИТОГОВЫЙ ЧЕК-ЛИСТ
% ============================================================

\section*{\color{darkaccent}10. Итоговый чек-лист}
\sectionline

Перед решением задачи проверьте, умеете ли Вы уверенно выполнить каждый пункт.

\begin{enumerate}[label=\(\square\)]
    \item Записать общую формулу
    \[
    x=x_0+v_xt.
    \]

    \item Составить отдельные формулы движения двух тел.

    \item Правильно определить знак скорости.

    \item Сформулировать условие встречи:
    \[
    x_1=x_2.
    \]

    \item Решить линейное уравнение относительно времени.

    \item Подставить найденное время в закон движения.

    \item Найти координату встречи.

    \item Проверить результат по формуле второго тела.

    \item Считать \(x_0\) с графика.

    \item Найти скорость по двум точкам:
    \[
    v_x=\frac{\Delta x}{\Delta t}.
    \]

    \item Отличить абсциссу от ординаты.

    \item Восстановить формулу движения по графику.

    \item Использовать известную координату встречи для обратных задач.
\end{enumerate}

\section*{\color{darkaccent}11. Формулы, которые нужно знать}
\sectionline

\[
\boxed{x=x_0+v_xt}
\]

\[
\boxed{x_1=x_{01}+v_{x1}t}
\]

\[
\boxed{x_2=x_{02}+v_{x2}t}
\]

\[
\boxed{x_1=x_2}
\qquad
\text{--- условие встречи}
\]

\[
\boxed{
v_x=\frac{\Delta x}{\Delta t}
=
\frac{x_2-x_1}{t_2-t_1}
}
\]

\[
\boxed{
x_{\text{встр}}
=
x_1(t_{\text{встр}})
=
x_2(t_{\text{встр}})
}
\]

% ============================================================
% ТРЕНИРОВОЧНЫЙ БЛОК
% ============================================================

\clearpage

\section*{\color{darkaccent}Тренировочный блок}
\sectionline

\begin{center}
{\large\textbf{Путь Б: 10.09--11.09}}
\end{center}

Все задачи относятся к равномерному прямолинейному движению вдоль оси \(Ox\).

\subsection*{\color{darkaccent}10.09.2026}

\begin{enumerate}
    \item Два тела движутся по законам
    \[
    x_1=2+3t,
    \qquad
    x_2=14-t.
    \]
    Найдите время и координату встречи.

    \item Начальные координаты двух тел:
    \[
    x_{01}=18\text{ м},
    \qquad
    x_{02}=-6\text{ м}.
    \]
    Скорости:
    \[
    v_{x1}=-2\text{ м/с},
    \qquad
    v_{x2}=4\text{ м/с}.
    \]
    Составьте законы движения и найдите время и координату встречи.

    \item График движения тела проходит через точки
    \[
    (0;7)
    \quad\text{и}\quad
    (4;3).
    \]
    Найдите \(x_0\), \(v_x\) и запишите закон движения.

    \item Первое тело движется по закону
    \[
    x_1=-4+2t.
    \]
    Второе тело движется по закону
    \[
    x_2=11-t.
    \]
    Найдите \(t_{\text{встр}}\) и \(x_{\text{встр}}\).

    \item Первое тело движется по закону
    \[
    x_1=1+2t.
    \]
    Через \(4\) с оно встречается со вторым телом, движение которого описывается формулой
    \[
    x_2=x_{02}+3t.
    \]
    Найдите координату встречи и \(x_{02}\).
\end{enumerate}

\vspace{6mm}

\subsection*{\color{darkaccent}11.09.2026}

\begin{enumerate}
    \item Два тела движутся по законам
    \[
    x_1=-5+4t,
    \qquad
    x_2=16-t.
    \]
    Найдите время и координату встречи.

    \item Начальные координаты:
    \[
    x_{01}=24\text{ м},
    \qquad
    x_{02}=-4\text{ м}.
    \]
    Скорости:
    \[
    v_{x1}=-3\text{ м/с},
    \qquad
    v_{x2}=4\text{ м/с}.
    \]
    Составьте законы движения и найдите точку встречи.

    \item График равномерного движения проходит через точки
    \[
    (0;-2)
    \quad\text{и}\quad
    (6;4).
    \]
    Найдите скорость и запишите формулу координаты.

    \item Для двух тел получены законы:
    \[
    x_1=8-\frac12t,
    \qquad
    x_2=-1+t.
    \]
    Найдите время и координату встречи.

    \item Первое тело движется по закону
    \[
    x_1=12-2t.
    \]
    В момент
    \[
    t_{\text{встр}}=3\text{ с}
    \]
    оно встречается со вторым телом:
    \[
    x_2=x_{02}+t.
    \]
    Найдите координату встречи и начальную координату второго тела.
\end{enumerate}

% ============================================================
% ОТВЕТЫ
% ============================================================

\clearpage

\section*{\color{darkaccent}Ответы к тренировочному блоку}
\sectionline

\subsection*{\color{darkaccent}10.09.2026}

\begin{center}
\renewcommand{\arraystretch}{1.45}
\begin{tabularx}{\linewidth}{@{}cX@{}}
\toprule
\textbf{№} & \textbf{Ответ} \\
\midrule
1 &
\(
t_{\text{встр}}=3\text{ с},
\quad
x_{\text{встр}}=11\text{ м}
\)
\\

2 &
\(
x_1=18-2t,
\quad
x_2=-6+4t,
\quad
t_{\text{встр}}=4\text{ с},
\quad
x_{\text{встр}}=10\text{ м}
\)
\\

3 &
\(
x_0=7\text{ м},
\quad
v_x=-1\text{ м/с},
\quad
x=7-t
\)
\\

4 &
\(
t_{\text{встр}}=5\text{ с},
\quad
x_{\text{встр}}=6\text{ м}
\)
\\

5 &
\(
x_{\text{встр}}=9\text{ м},
\quad
x_{02}=-3\text{ м}
\)
\\
\bottomrule
\end{tabularx}
\end{center}

\vspace{8mm}

\subsection*{\color{darkaccent}11.09.2026}

\begin{center}
\renewcommand{\arraystretch}{1.45}
\begin{tabularx}{\linewidth}{@{}cX@{}}
\toprule
\textbf{№} & \textbf{Ответ} \\
\midrule
1 &
\(
t_{\text{встр}}=\frac{21}{5}\text{ с}=4{,}2\text{ с},
\quad
x_{\text{встр}}=\frac{59}{5}\text{ м}=11{,}8\text{ м}
\)
\\

2 &
\(
x_1=24-3t,
\quad
x_2=-4+4t,
\quad
t_{\text{встр}}=4\text{ с},
\quad
x_{\text{встр}}=12\text{ м}
\)
\\

3 &
\(
v_x=1\text{ м/с},
\quad
x=-2+t
\)
\\

4 &
\(
t_{\text{встр}}=6\text{ с},
\quad
x_{\text{встр}}=5\text{ м}
\)
\\

5 &
\(
x_{\text{встр}}=6\text{ м},
\quad
x_{02}=3\text{ м}
\)
\\
\bottomrule
\end{tabularx}
\end{center}

\end{document}